\documentclass[a4paper,10pt]{article}

\usepackage[utf8]{inputenc}
\usepackage[T1]{fontenc}
\usepackage[ngerman]{babel}
\usepackage{amsmath}
\usepackage{amssymb}
\usepackage{enumitem}
\usepackage{xcolor}
\usepackage{geometry}
\geometry{a4paper,margin=1.9cm}

% ----- Dark Mode (kitty-Theme: kitty.conf) -----
\definecolor{bgdark}{HTML}{11121A}
\definecolor{fgtext}{HTML}{D5DDE8}
\definecolor{accent}{HTML}{91A8D0}   % blau (Kapitel)
\definecolor{accent2}{HTML}{E8D49A}  % gelb (Aufgaben)
\definecolor{good}{HTML}{B5D4A8}     % grün (Ergebnis/Lösung)
\definecolor{rulecol}{HTML}{4A5B78}
\pagecolor{bgdark}
\color{fgtext}

\newcommand{\kapitel}[1]{%
  \par\vspace{6pt}
  {\large\bfseries\color{accent} #1}\par
  \noindent\textcolor{accent}{\rule{\linewidth}{0.8pt}}\par\vspace{2pt}}
\newcommand{\aufg}[2]{%
  \par\vspace{4pt}
  {\bfseries\color{accent2} #1}~~{\footnotesize\color{rulecol}[#2]}\par\vspace{1pt}}
\newcommand{\loes}{\textbf{\color{good}Lösung.}~}
\newcommand{\erg}[1]{\textcolor{good}{\textbf{#1}}}

\setlist[itemize]{leftmargin=1.3em,itemsep=1pt,topsep=1pt,parsep=0pt,label=\textcolor{good}{\textbullet}}
\setlist[enumerate]{leftmargin=1.8em,itemsep=1pt,topsep=1pt,parsep=0pt}
\setlength{\parindent}{0pt}
\setlength{\parskip}{2.5pt}

\begin{document}

\begin{center}
  {\Large\bfseries\color{accent} Numerische Methoden -- Musterlösungen}\\[1pt]
  {\small Vollständige Rechenwege zu \texttt{aufgaben\_dark} / \texttt{aufgaben\_white}
   -- gleiche Nummerierung, alle Werte eigenständig nachgerechnet.}
\end{center}
\noindent\textcolor{rulecol}{\rule{\linewidth}{0.4pt}}

{\footnotesize Hinweis: Einige durchgerechnete Beispiele im Skript enthalten
Übertragungsfehler (z.\,B.\ in der Newton-/Sekanten-Iteration und bei Mittelpunkt-/
Trapez-Werten); hier stehen die \emph{korrekten} Werte nach derselben Methode.}

% =====================================================================
\kapitel{Kapitel 1 \;--\; Einstieg in die Numerischen Methoden}

\aufg{1.1 \;--\; Minimum der Sinusfunktion}{Handrechnung}
\loes Exakt: $f'(x)=\cos x=0\Rightarrow x=-\tfrac{\pi}{2}\approx-1{,}571$, also
$\erg{f=-1}$ (Ränder $\sin(-3)\approx-0{,}141$, $\sin(1)\approx0{,}841$). Gitter
$h=1$ (Knoten $-3,-2,-1,0,1$): kleinster Wert $\sin(-2)=\erg{-0{,}909}$ bei $x=-2$.
Fehler $|-1-(-0{,}909)|=0{,}091$ -- das grobe Gitter trifft das wahre Minimum bei
$-\tfrac\pi2$ nicht.

\aufg{1.2 \;--\; Grid-Search}{Handrechnung}
\loes Über alle $5\times5=25$ Paare $|2w+b-11|+|5w+b-13|$ auswerten; Minimum bei
$(w,b)=(0,10)$: $|10-11|+|10-13|=1+3=\erg{4}$. Exakt wäre $w=\tfrac23$,
$b=\tfrac{29}{3}\approx9{,}667$; das Gitter trifft nur näherungsweise. Aufwand: $25$
Auswertungen, wächst exponentiell mit der Variablenzahl ($k^d$).

\aufg{1.3 \;--\; Enter the machine}{Code}
\loes Doppelschleife über das $(w,b)$-Gitter, je Paar den Fehler berechnen, Minimum
mitführen $\to(0,10)$. Zeigt die $O(k^d)$-Kosten der Brute-Force.

\aufg{Denkanstoß / Selbstreflexion}{kurz}
\loes Verfeinerung: erst grobes Gitter, dann lokal verfeinern (coarse-to-fine) oder
Gradienteninformation nutzen. Feineres $h$ $\Rightarrow$ genauer, aber mehr
Auswertungen; $[a,b]$ muss das Extremum enthalten; Diskretisierung ersetzt das
Kontinuum durch endlich viele Punkte $\Rightarrow$ exakte Extrema fallen dazwischen.

% =====================================================================
\kapitel{Kapitel 2 \;--\; Gleitkomma-Arithmetik}

\aufg{2.1 \;--\; Rundungsfehler}{Handrechnung}
\loes $10/3=3{,}\overline{3}$; nach drei Stellen $3{,}333$. Fehler
$|3{,}3333\ldots-3{,}333|=\erg{$0{,}000\overline{3}=\tfrac{1}{3000}$}$.

\aufg{2.2 \;--\; 2-Bit-Mantisse \& Maschinenepsilon}{Handrechnung}
\loes Signifikand $0{.}b_1b_2$: $\{0;\,0{,}25;\,0{,}5;\,0{,}75\}$;
$\varepsilon_{\mathrm{mach}}=2^{-2}=\erg{0{,}25}$. Relative Präzision $\approx2^{-t}$
konstant; absolute Präzision $\varepsilon_{\mathrm{mach}}\cdot|x|$ wächst mit dem
Betrag.

\aufg{2.3 \;--\; Auslöschung von Hand}{Handrechnung}
\loes \textbf{(a)} $\tilde a=12345679$, $\tilde b=12345678\Rightarrow\tilde a-\tilde b=1$,
aber $a-b=0{,}5$, Fehler $\erg{0{,}5}$. \textbf{(b)} $\tilde a=\tilde b=12345678\Rightarrow0$,
aber $a-b=0{,}4$. Wahre Differenzen liegen nah beieinander, das Maschinenergebnis
springt $0\!\leftrightarrow\!1$ -- riesiger relativer Fehler (katastrophale
Auslöschung).

\aufg{2.4 \;--\; Festkomma-Skalierung}{Handrechnung}
\loes $1{,}234567\cdot10^{6}=\erg{1234567}$; Präzision $10^{-6}=0{,}000001$; maximaler
Rundungsfehler $0{,}5\cdot10^{-6}=\erg{0{,}0000005}$. Über den ganzen Bereich konstant.

\aufg{2.5 \;--\; Maschinenepsilon (hands on)}{Code}
\loes $\varepsilon\leftarrow1$; solange $1+\varepsilon/2\neq1$: $\varepsilon\leftarrow\varepsilon/2$.
Ergebnis $\approx2^{-t}$ (float32 $\approx1{,}19\cdot10^{-7}$, double
$\approx2{,}22\cdot10^{-16}$). Festkomma: konstanter kleinster Abstand; Gleitkomma:
relativ zur Größe. Taschenrechner: \erg{double} (genug relative Präzision, großer Bereich).

\aufg{2.6 \;--\; Auslöschung demonstrieren}{Code}
\loes $g(\varepsilon)=1/(1+\varepsilon-1)$ vs.\ $1/\varepsilon$ auswerten: für große
$\varepsilon$ gleich, für sehr kleine dominiert der Rundungsfehler (Nenner-Auslöschung),
das Ergebnis explodiert.

\aufg{2.7 \;--\; Argumentation}{Konzept}
\loes Für $a=b$ ist $f(\epsilon)=1/(a+\epsilon-b)=1/\epsilon$ -- für kleines
$\epsilon$ \erg{explodiert} der Wert, und der Nenner $a+\epsilon-b$ leidet unter
Auslöschung (führende Stellen heben sich auf). \emph{Randnotiz $a>b$:} dann dominiert
$a-b$, keine Auslöschung, $f$ bleibt stabil.

\aufg{Denkanstoß / Selbstreflexion}{kurz}
\loes Fehlerbasierte Metriken $\to$ Kondition/Stabilität/Konsistenz/Konvergenz (Kap.\,3).
Gleitkomma $x=(-1)^s m\,2^e$ (Vorzeichen, Mantisse=Präzision, Exponent=Größe);
$\varepsilon_{\mathrm{mach}}=2^{-t}$ ist die obere Schranke des relativen
Rundungsfehlers; praktisch akkumulieren Rundungsfehler in Iterationen.

% =====================================================================
\kapitel{Kapitel 3 \;--\; Fehleranalyse}

\aufg{3.1 \;--\; Minimum der Sinusfunktion (Eigenschaften)}{Handrechnung}
\loes \textbf{Kondition} $\kappa=|f(x)-f(\tilde x)|$ (problemabhängig):
\[
  \kappa(-1,\!+0{,}25)=|\sin(-1)-\sin(-0{,}75)|=0{,}1599,\quad
  \kappa(-1,\!-0{,}25)=0{,}1074,\quad \kappa(0,\pm0{,}25)=0{,}2474.
\]
Um $x=-1$ \erg{gut konditioniert} ($0{,}11$--$0{,}16$), um $x=0$ schlechter ($0{,}25$).\\
\textbf{Stabilität} ($\tilde x=x+0{,}25$, $x=-1$): $h=1$ -- beide auf demselben
Knoten $\Rightarrow s\approx\erg{0}$; $h=0{,}2$ --
$s=\frac{|\sin(-0{,}8)-\sin(-1)|}{0{,}25}=\frac{0{,}1241}{0{,}25}\approx\erg{0{,}496}$.\\
\textbf{Konsistenz} $c\ge|\hat f(x)-f(x)|$, wahres Minimum $f(-\tfrac\pi2)=-1$:
$c_{h=1}\ge|-0{,}841-(-1)|=\erg{0{,}159}$,
$c_{h=0{,}2}\ge|-0{,}949-(-1)|=\erg{0{,}051}$ -- verbessert sich mit kleinerem $h$.\\
\textbf{Konvergenz} $|\hat f(\tilde x)-f(x)|\approx0{,}16$ in beiden Fällen; bei
$h=0{,}2$ \emph{keine} Verbesserung, weil die geringere Stabilität die Störung
aufbläht. Merke: Konvergenz $=$ Stabilität $+$ Konsistenz.

\aufg{3.2 \;--\; Stabilität und Kondition}{Beweis}
\loes Für $h\to0$ wird das Verfahren exakt, $\hat f\to f$, also
$|\hat f(\tilde x)-\hat f(x)|\to|f(\tilde x)-f(x)|=\kappa$. Die rechte Seite ist die
Kondition: bei perfekter Diskretisierung erbt das Verfahren genau die
Problem-Empfindlichkeit. $\square$

\aufg{3.3 \;--\; Compute-getriebene Fehleranalyse}{Code}
\loes Automatisiert über das $(w,b)$-System: Schrittweite/Störung variieren, $\kappa$,
$s$, $c$ und Gesamtfehler messen und über die Methode optimieren.

\aufg{Selbstreflexion / Denkanstoß}{kurz}
\loes Gute Konsistenz (kleiner Bias) $+$ gute Stabilität (kleine Varianz) $=$ gute
Konvergenz. Für $h\to0$ Grenze durch Maschinengenauigkeit (Rundungsrauschen) $\Rightarrow$
es gibt ein optimales $h$. Verfeinerung von Brute-Force: coarse-to-fine.

% =====================================================================
\kapitel{Kapitel 4 \;--\; Newton-Verfahren}

\aufg{4.1 \;--\; Newton für $\sqrt2$}{Handrechnung}
\loes $f(x)=x^2-2$, $f'(x)=2x$, also
$x_{n+1}=x_n-\frac{x_n^2-2}{2x_n}=\tfrac12\!\left(x_n+\tfrac{2}{x_n}\right)$, $x_0=2$:
\[
\begin{aligned}
x_1&=\tfrac12\!\left(2+\tfrac22\right)=\tfrac12(2+1)=1{,}5,\\
x_2&=\tfrac12\!\left(1{,}5+\tfrac{2}{1{,}5}\right)=\tfrac12(1{,}5+1{,}33333)=1{,}41667,\\
x_3&=\tfrac12\!\left(1{,}41667+\tfrac{2}{1{,}41667}\right)=\tfrac12(1{,}41667+1{,}41176)=1{,}414216,\\
x_4&=\tfrac12\!\left(1{,}414216+\tfrac{2}{1{,}414216}\right)=\erg{1{,}41421356}.
\end{aligned}
\]
Fehler: $0{,}59\to0{,}086\to0{,}0025\to2{,}1\cdot10^{-6}$ -- quadriert sich pro Schritt.

\aufg{4.2 \;--\; Newton vs.\ Grid-Search für $\sin(x)=0$}{Handrechnung}
\loes \textbf{(a)} Grid $h=0{,}3$: $f(3{,}1)\approx0{,}042$ am nächsten an $0$ ($6$
Auswertungen), Fehler $|3{,}1-\pi|\approx0{,}042$. \textbf{(b)} Newton $f'=\cos$,
$x_0=3{,}5$: $x_1=3{,}5-\frac{-0{,}351}{-0{,}936}=3{,}125$, $x_2\approx\erg{3{,}14159}$.
\textbf{(c)} Newton: Fehler $<10^{-5}$ in $2$ Schritten; Grid nur $0{,}042$ -- die
Ableitung lenkt effizienter.

\aufg{4.3 \;--\; Versagensmodi}{Schaubild}
\loes Drei Modi: (i) waagrechte Tangente $f'(x_0)=0$ $\Rightarrow$ Schritt nach
$\pm\infty$; (ii) Oszillation/Zyklus bei symmetrischen Punkten; (iii) mehrdeutige
Startpunktwahl $\Rightarrow$ Konvergenz zur falschen Wurzel.

\aufg{4.4 \;--\; Newton von Hand}{Handrechnung}
\loes $f=x^3-x$, $f'=3x^2-1$, $x_{n+1}=x_n-\frac{x_n^3-x_n}{3x_n^2-1}$, $x_0=1{,}4$:
\[
\begin{aligned}
x_1&=1{,}4-\frac{1{,}4^3-1{,}4}{3\cdot1{,}4^2-1}=1{,}4-\frac{2{,}744-1{,}4}{5{,}88-1}=1{,}4-\frac{1{,}344}{4{,}88}=1{,}1246,\\
x_2&=1{,}1246-\frac{1{,}1246^3-1{,}1246}{3\cdot1{,}1246^2-1}=1{,}1246-\frac{0{,}2977}{2{,}7941}=1{,}0180,\\
x_3&=1{,}0180-\frac{1{,}0180^3-1{,}0180}{3\cdot1{,}0180^2-1}=1{,}0180-\frac{0{,}0371}{2{,}1092}=1{,}0005,\\
x_4&\approx\erg{1{,}0000}.
\end{aligned}
\]
Quadratische Konvergenz: die Zahl korrekter Stellen verdoppelt sich pro Schritt.

\aufg{4.5 \;--\; Sekantenverfahren von Hand}{Handrechnung}
\loes $x_{n+1}=x_n-f(x_n)\frac{x_n-x_{n-1}}{f(x_n)-f(x_{n-1})}$, Startwerte
$f(1{,}2)=0{,}528$, $f(1{,}4)=1{,}344$:
\[
\begin{aligned}
x_2&=1{,}4-1{,}344\cdot\frac{1{,}4-1{,}2}{1{,}344-0{,}528}=1{,}4-1{,}344\cdot\frac{0{,}2}{0{,}816}=1{,}0706,\\
x_3&=1{,}0706-0{,}1565\cdot\frac{1{,}0706-1{,}4}{0{,}1565-1{,}344}=1{,}0706-0{,}1565\cdot\frac{-0{,}3294}{-1{,}1875}=1{,}0272,\\
x_4&=1{,}0272-0{,}0566\cdot\frac{1{,}0272-1{,}0706}{0{,}0566-0{,}1565}=1{,}0272-0{,}0246=\erg{1{,}0026}.
\end{aligned}
\]
mit $f(1{,}0706)=0{,}1565$, $f(1{,}0272)=0{,}0566$. Superlinear (Ordnung
$\approx1{,}618$), ohne Ableitung.

\aufg{4.6 \;--\; Geometrische Fehlersuche}{Handrechnung}
\loes Newton-Abb.\ $N(x)=\frac{2x^3}{3x^2-1}$.
(i) $f'(x_0)=0\Rightarrow x_0=\pm\tfrac{1}{\sqrt3}\approx\erg{$\pm0{,}577$}$ (Tangente
waagrecht). (ii) falsche Wurzel: $x_0=0{,}4\to$ Becken von $x=0$. (iii) Oszillation:
$N(x_0)=-x_0\Rightarrow 5x_0^2=1\Rightarrow x_0=\tfrac{1}{\sqrt5}\approx\erg{0{,}447}$
(Zyklus $\pm0{,}447$).

\aufg{4.7 \;--\; Enter the machine}{Code}
\loes Drei Solver auf $x^3-x$: Grid linear, Newton quadratisch, Sekanten superlinear;
Startpunkte $\{\pm\tfrac1{\sqrt3},\tfrac1{\sqrt5}\}$ reproduzieren die Versagensmodi.

\aufg{Selbstreflexion / Denkanstoß}{kurz}
\loes Newton nutzt Steigung (Taylor), konvergiert quadratisch: $10^{-1}\to10^{-16}$
in $\sim4$ Schritten (Stellenverdopplung). Sekanten, wenn $f'$ teuer/unbekannt -- nicht
bei fehlenden guten Startwerten. Global finden $\to$ Kap.\,5.

% =====================================================================
\kapitel{Kapitel 5 \;--\; Globale Optimierung}

\aufg{5.1 \;--\; Newton auf Shekel}{Code}
\loes Newton/Gradientenabstieg konvergieren zum \emph{nächsten} Foxhole $\Rightarrow$
Ergebnis hängt von $x_0$ ab. \emph{Randnotiz Maxima:} Vorzeichen umkehren,
$x_{n+1}=x_n+\frac{f'(x_n)}{|f''(x_n)|}$ (Schritt bergauf).

\aufg{5.2 \;--\; Random Restarts}{Handrechnung}
\loes \textbf{(a)} $P=1-0{,}85^{10}=\erg{0{,}803}$ ($\approx80\%$).
\textbf{(b)} $K=\frac{\ln0{,}01}{\ln0{,}85}=28{,}3\Rightarrow\erg{29}$.
\textbf{(c)} $p=0{,}05:\,K=\frac{-4{,}605}{\ln0{,}95}=89{,}8\Rightarrow\erg{90}$;
$p=0{,}01:\,K=458{,}2\Rightarrow\erg{459}$. Kleineres Becken $\Rightarrow$ drastisch mehr
Neustarts.

\aufg{5.3--5.10 \;--\; Code \& Konzept}{Modellantworten}
\loes
\begin{itemize}
  \item \textbf{5.3} Shekel als Summe negativer Foxholes; Newton bleibt im nächsten
        Becken / versagt bei $f'=0$.
  \item \textbf{5.4} $|f''|$-Trick: Krümmungs-Vorzeichen erzwingen $\Rightarrow$ Schritt
        stets bergab; geometrisch wird die Tangente gespiegelt, projiziert nie bergauf.
  \item \textbf{5.5} Größtes Einzugsgebiet gehört meist zum globalen Minimum; seine
        relative Größe ist gerade das $p$ aus 5.2.
  \item \textbf{5.6 / 5.10} Basin Hopping (lokale Sprünge) braucht meist \emph{weniger}
        Auswertungen als Random Restarts -- nutzt Nähe; zu kleine Sprünge $\to$ stecken,
        zu große $\to$ reine Zufallssuche.
  \item \textbf{5.7} Schwer abgrenzbar bei hochfrequenten Funktionen, z.\,B.\
        $g(x)=\sin(x)+0{,}1\sin(50x)$.
  \item \textbf{5.8} Heuristik: Sprungweite bei Stagnation vergrößern, nach Erfolg
        verkleinern; katastrophale Sprünge via Clipping / Metropolis-Akzeptanz verhindern.
  \item \textbf{5.9} Mini-Batches liefern eine verrauschte Loss-Landschaft; das Rauschen
        schüttelt den Optimierer aus flachen lokalen Minima (Glättung).
\end{itemize}

\aufg{5.11 \;--\; Einzugsgebiete kartieren}{Handrechnung}
\loes \textbf{(a)} Grenzen mittig: $\tfrac{0+4}{2}=2$, $\tfrac{4+7}{2}=5{,}5$, also
$(-\infty,2)\!\to\!0$, $(2;5{,}5)\!\to\!4$, $(5{,}5,\infty)\!\to\!7$.
\textbf{(b)} Global bei $4$, Becken $(2;5{,}5)$. Sprung ab $7$: $7+\delta\in(4,10)$;
Treffer wenn $7+\delta<5{,}5\Leftrightarrow\delta\in(-3;-1{,}5)$, Länge $1{,}5$:
$P=\frac{1{,}5}{6}=\erg{0{,}25}$.
\textbf{(c)} Basin Hopping springt aus $x=7$ ins \emph{benachbarte} globale Becken bei
$4$ (Nähe genutzt); Random Restarts würfeln über den ganzen Raum, in dem das globale
Becken nur klein ist.

\aufg{Selbstreflexion / Denkanstoß}{kurz}
\loes Lokale Verfahren folgen nur lokaler Information $\to$ nächstes Becken. Aufwand
$\sim1/p$. $f''(x_n)$ misst Krümmung (Beckenbreite, min/max/Sattel). Hochfrequente
Shekel-Version $\to$ viele winzige verschachtelte Becken, extrem schwer.

% =====================================================================
\kapitel{Kapitel 6 \;--\; Numerische Integration}

\noindent Referenz: $\tilde I=\int_{-2}^{-1}\sin x\,dx=[-\cos x]_{-2}^{-1}
=-\cos(-1)+\cos(-2)=\erg{-0{,}95645}$. Werte: $\sin(-2)=-0{,}90930$,
$\sin(-1{,}75)=-0{,}98399$, $\sin(-1{,}5)=-0{,}99750$, $\sin(-1{,}25)=-0{,}94898$,
$\sin(-1)=-0{,}84147$.

\aufg{6.1 \;--\; Exakter Wert}{Handrechnung}
\loes Siehe oben: $\erg{-0{,}95645}$.

\aufg{6.2 \;--\; Mittelpunktsregel, $n=1$}{Handrechnung}
\loes $m=-1{,}5$: $\hat I=1\cdot\sin(-1{,}5)=\erg{-0{,}99750}$, Fehler $0{,}04105$.

\aufg{6.3 \;--\; Mittelpunktsregel, $n=2$}{Handrechnung}
\loes $h=0{,}5$: $\hat I=0{,}5[\sin(-1{,}75)+\sin(-1{,}25)]=\erg{-0{,}96649}$, Fehler
$0{,}01004$ (viertelt sich gegenüber $n=1$).

\aufg{6.4 \;--\; Trapezregel, $n=2$}{Handrechnung}
\loes $T_2=\tfrac{h}{2}[f(-2)+2f(-1{,}5)+f(-1)]=0{,}25[-0{,}90930-1{,}99499-0{,}84147]=\erg{-0{,}93644}$,
Fehler $0{,}02001$.

\aufg{6.5 \;--\; Simpson-Regel, $n=2$}{Handrechnung}
\loes $S_2=\tfrac{h}{3}[f(-2)+4f(-1{,}5)+f(-1)]=\tfrac{0{,}5}{3}[-0{,}90930-3{,}98998-0{,}84147]=\erg{-0{,}95679}$,
Fehler $0{,}00034$ (Ordnung $4$, weit genauer; nicht exakt, da $\sin$ kein Kubik ist).

\aufg{6.6 \;--\; Monte Carlo, $N=4$}{Handrechnung}
\loes $f$-Werte $-0{,}92896,-0{,}99978,-0{,}91276,-0{,}86742$;
$\hat I=1\cdot\frac{-3{,}70893}{4}=\erg{-0{,}92723}$, Fehler $0{,}02922$.

\aufg{6.7 \;--\; Monte-Carlo-Fehler}{Konzept}
\loes $\bar f=-0{,}92723$,
$\sigma_f^2=\frac{1}{3}\sum(f-\bar f)^2\approx\erg{0{,}00302}$, $\sigma_f\approx0{,}0549$,
$\mathrm{SE}=\frac{(b-a)\sigma_f}{\sqrt N}=\frac{0{,}0549}{2}\approx\erg{0{,}0275}$. Bei
$N=16$: $\mathrm{SE}\propto1/\sqrt N\Rightarrow$ halbiert auf $\approx0{,}0137$.
Dimensionsunabhängig ($d$ taucht nicht auf).

\aufg{6.8 \;--\; Simpson-Gewichte via Lagrange}{Herleitung}
\loes Mit $x_0=a,x_1=m,x_2=b$ und Lagrange-Basis $L_a,L_m,L_b$; Integration über
$[a,b]$ liefert die Gewichte $\frac{b-a}{6},\frac{4(b-a)}{6},\frac{b-a}{6}$, also
\[
  \int_a^b f\,dx\approx\frac{b-a}{6}\big[f(a)+4f(m)+f(b)\big].
\]
Da Grad-$\le2$-Polynome exakt interpoliert werden, ist Simpson für quadratische (und
kubische) Funktionen exakt. $\square$

\aufg{6.9 \;--\; Höhere Dimensionen}{Code}
\loes $N$ Zufallspunkte im Hyperwürfel, $f$ mitteln, mit Volumen multiplizieren; Fehler
$\propto1/\sqrt N$ \erg{unabhängig von $d$}, während Gitterquadratur $k^d$ Punkte
braucht (Fluch der Dimensionalität).

\aufg{Selbstreflexion / Denkanstoß}{kurz}
\loes Rangfolge: Simpson ($10^{-4}$) $\gg$ Mittelpunkt ($10^{-2}$) $>$ Trapez
$\approx$ MC. Mittelpunkt schlägt Endpunktregeln (lineare Fehler heben sich auf).
Deterministische Quadratur scheitert in hohen Dim.; MC nicht ($1/\sqrt N$). Mittelpunkt
$=$ MC mit mittigen, nicht-zufälligen Punkten. Mini-Batch-Mittel in SGD $=$ MC-Schätzer
des Gradienten. Hochgradige Quadratur wird instabil (Runge-artige Oszillation).

% =====================================================================
\kapitel{Kapitel 7 \;--\; Modelltraining}

\aufg{7.1 \;--\; Loss bei Initialisierung}{Handrechnung}
\loes Mit $w=0,b=\bar y$ ist die Vorhersage konstant $\bar y$, also
$L(0,\bar y)=\tfrac1n\sum(\bar y-y_i)^2=\mathrm{Var}(y)=\frac{1401{,}5}{14}\approx\erg{100{,}1}$.
$5000\approx$ summierter statt gemittelter Loss; $0{,}0001\approx$ ungewollte
Zusatz-Normierung.

\aufg{7.2 \;--\; Overfitting als Sanity-Check}{Konzept}
\loes Durch zwei Punkte gehen \erg{unendlich viele} Grad-4-Kurven (5 Parameter, 2
Bedingungen $\Rightarrow$ 3-Parameter-Schar) $\Rightarrow$ unendlich viele globale
Optima, hochgradig entartete Verlustlandschaft.

\aufg{7.3 \;--\; Mini-Batch-SGD}{Konzept}
\loes Mini-Batch-SGD mittelt Verlust-/Gradienten-Beiträge über eine zufällige
Daten-Stichprobe $\Rightarrow$ Monte-Carlo-Schätzer der Erwartung (Standardfehler
$\propto1/\sqrt B$). Der Engpass kippt, wenn der Voll-Batch nicht mehr in den Speicher
passt bzw.\ seine Kosten den Varianzvorteil übersteigen (großes $n$).

\aufg{7.4 \;--\; Mehrere Random Seeds}{Konzept}
\loes Für $\hat y=w^2x+b$ ist $\partial L/\partial w\propto w$, am Ursprung $=0$ bei
Krümmung in $b$ und Flachheit in $w$ $\Rightarrow$ \erg{Sattel}; $w=0$ bleibt hängen.
Mehrere Seeds $=$ Random Restarts auf der nicht-konvexen Landschaft (verschiedene Becken).

\aufg{7.5 \;--\; Quantisierung für Inferenz}{Konzept}
\loes Inferenz akkumuliert keine Updates $\Rightarrow$ toleriert grobe Präzision;
Training braucht Präzision für die kleinen Gradientenschritte. float32 wird versendet
für Reserve an Wertebereich/Genauigkeit (und teils aus Trägheit).

\aufg{Selbstreflexion}{kurz}
\loes Loss$\to0$ bei Überparametrisierung garantiert erreichbar $\Rightarrow$ Misserfolg
= Pipeline-Bug, nicht Optimierer. Erster Loss $\approx100=\mathrm{Var}(y)$; $5000$
summiert, $0{,}0001$ über-normiert. SGD-Mini-Batch $=$ MC des Gradienten. $w^2x+b$:
Ursprung Sattel, Null-Init Falle. Hoch trainieren / niedrig einsetzen funktioniert, weil
Inferenz keine Updates akkumuliert.

\end{document}
